%% toc: P3 31 31.5,31.6,31.7,31.8
\chapter{Two-Cardinal Models and Cardinal Transfer}
\label{sec:p03-w03-p03-029-two-cardinal}

A first-order theory with a distinguished unary predicate can control two
sizes at once: the size of a model and the size of the part of the model
picked out by the predicate. A \emph{two-cardinal theorem} says that if a
theory realizes one pair of such sizes, then it realizes another. Several
results of this shape carry similar names and are easily confused. They
differ in their hypotheses, in the pairs they reach and in their methods. This
chapter fixes the convention for pairs of cardinals, states Vaught's
\((\aleph_1,\aleph_0)\) theorem in the form in which it was first proved, and
separates it from the L\"owenheim--Skolem theorems and from the far-apart and
gap-one transfer theorems.

\emph{Inputs.} A countable first-order language with a distinguished unary
predicate, a theory in that language, and a model in which the predicate is
infinite and strictly smaller than the universe. \emph{Conclusion.} A model
of the same complete theory whose universe has cardinality \(\aleph_1\) and
whose predicate is countably infinite. The downward L\"owenheim--Skolem
theorem (\cref{thm:downward-ls}) and the Ehrenfeucht--Mostowski construction
(\cref{thm:p03-em-models}) are used only where named. The statement-level
entry \cref{thm:p03-vaught-two-cardinal} is the same theorem as
\cref{thm:p03-w03-p03-029-theory-form} below.

\section{Vaught's theorem and distinct transfer theorems}
\label{sec:p03-w03-p03-029-vaught}

\subsection{Pairs of cardinals}

\begin{hbdefinition}{def:p03-w03-p03-029-kappa-lambda-model}{\((\kappa,\lambda)\)-models and two-cardinal models}
Let \(\Lang\) be a first-order language and let \(U\) be a unary relation
symbol of \(\Lang\), called the \emph{distinguished predicate}. For an
\(\Lang\)-structure \(\mathcal M\) write \(U^{\mathcal M}\subseteq\dom\mathcal
M\) for the interpretation of \(U\). Let \(\kappa\) and \(\lambda\) be
cardinals.
\begin{enumerate}
  \item \(\mathcal M\) is a \emph{\((\kappa,\lambda)\)-model}, or a model
        \emph{of type} \((\kappa,\lambda)\), if
        \(\abs{\dom\mathcal M}=\kappa\) and \(\abs{U^{\mathcal M}}=\lambda\).
        The first coordinate always measures the universe and the second
        always measures the distinguished predicate.
  \item An \(\Lang\)-theory \(T\) \emph{admits} \((\kappa,\lambda)\) if some
        model of \(T\) is a \((\kappa,\lambda)\)-model.
  \item A \((\kappa,\lambda)\)-model is a \emph{two-cardinal model} if
        \(\kappa>\lambda\geq\aleph_0\): the predicate is infinite and the
        universe is strictly larger than it.
\end{enumerate}
Since \(U^{\mathcal M}\subseteq\dom\mathcal M\), every
\((\kappa,\lambda)\)-model has \(\lambda\leq\kappa\). The definition allows a
finite \(\lambda\) and allows \(\lambda=\kappa\); neither case is a
two-cardinal model. The choice of \(U\) is part of the data: the same
structure can be a two-cardinal model for one unary symbol and fail to be one
for another.

This is the convention used for every two-cardinal statement in this
handbook. In \cref{thm:p03-vaught-two-cardinal} the hypothesis
\(\abs M>\abs{U^M}\geq\aleph_0\) says exactly that \(M\) is a two-cardinal
model for \(U\), and the conclusion says that \(T\) admits
\((\aleph_1,\aleph_0)\).

Morley and Vaught write \(\overline{\mathfrak A}\) for the power
(cardinality) of a relational system \(\mathfrak A\)
\cite[p.~38]{MorleyVaught1962}, and their two-cardinal theorem concerns a
singulary relation \(R_j^{\mathfrak A}\) with
\(\omega\leq\overline{R_j^{\mathfrak A}}<\overline{\mathfrak A}\)
\cite[p.~55, Theorem~6.2]{MorleyVaught1962}; there \(R_j\) plays the role of
\(U\), and \(\omega\), \(\omega_0\), \(\omega_1\) are the cardinals
\(\aleph_0\), \(\aleph_0\), \(\aleph_1\). Shelah works in a countable
vocabulary \(\tau\) with a distinguished unary relation symbol \(P\) and writes
\[
  K_{(\lambda,\mu)}=\{M:\lVert M\rVert=\lambda\ \&\ \abs{P^M}=\mu\}
\]
\cite[p.~3, 0.1~Notation]{Shelah2004PairAlephN}. His letters differ from the
ones used here, but the coordinate order is the same: \(K_{(\lambda,\mu)}\) is
the class of \((\lambda,\mu)\)-models in the sense of item~1, with
\(\lVert M\rVert\) the cardinality of the universe of \(M\).
\index{two-cardinal model}\index{model!of type \((\kappa,\lambda)\)}
\end{hbdefinition}

\begin{hbconvention}{def:p03-w03-p03-029-transfer-arrow}{Transfer arrows between pairs}
Shelah's notation records a transfer between pairs
\cite[p.~3, 0.2~Definition]{Shelah2004PairAlephN}. For infinite cardinals
\(\lambda,\mu,\lambda',\mu',\kappa\):
\begin{enumerate}
  \item \((\lambda,\mu)\to_\kappa(\lambda',\mu')\) means that every
        first-order theory \(T\) in \(\mathbb L(\tau)\) with \(\abs T<\kappa\)
        that has a model in \(K_{(\lambda,\mu)}\) also has a model in
        \(K_{(\lambda',\mu')}\);
  \item \((\lambda,\mu)\to'_\kappa(\lambda',\mu')\) means that every such
        \(T\) with \(\abs T<\kappa\) has a model in \(K_{(\lambda',\mu')}\)
        provided every finite subset of \(T\) has a model in
        \(K_{(\lambda,\mu)}\);
  \item the subscript \(\kappa^+\) may be written \({\leq}\kappa\), and the
        subscript is omitted when \(\kappa=\aleph_0\).
\end{enumerate}
The hypothesis of \(\to'\) is weaker than that of \(\to\), so
\((\lambda,\mu)\to'_\kappa(\lambda',\mu')\) implies
\((\lambda,\mu)\to_\kappa(\lambda',\mu')\). With \(\kappa=\aleph_1\) the
theories in question are the countable ones. An arrow states a transfer for a
fixed source pair and a fixed target pair; it says nothing about the pairs
that the theory does not realize.
\end{hbconvention}

\subsection{The theorem in its source form}

The \((\aleph_1,\aleph_0)\) theorem was announced by Vaught in 1961 and
proved by Morley and Vaught as Theorem~6.2 of \emph{Homogeneous universal
models}, in a section titled a L\"owenheim--Skolem theorem for two cardinals
\cite[p.~55 and reference~25 on p.~57]{MorleyVaught1962}. From p.~54 on,
the paper restricts itself to relational systems with countable index sets,
that is, to countably many relation symbols and no function symbols
\cite[p.~54]{MorleyVaught1962}.

\begin{hbtheorem}{thm:p03-w03-p03-029-mv-6-2}{proved=annex:p03-w03-p03-029-mv-6-2}{Morley--Vaught Theorem 6.2}
Let \(\mathfrak A\) be a relational system with a countable index set. If
\(R_j^{\mathfrak A}\) is singulary and
\(\omega\leq\overline{R_j^{\mathfrak A}}<\overline{\mathfrak A}\), then
there exists a system \(\mathfrak B\equiv\mathfrak A\) such that
\(\overline{\mathfrak B}=\omega_1\) and
\(\overline{R_j^{\mathfrak B}}=\omega_0\)
\cite[p.~55, Theorem~6.2; standing hypothesis p.~54]{MorleyVaught1962}.
The proof given here uses the Axiom of Choice through the
L\"owenheim--Skolem theorem and the choices made along the recursion.
\end{hbtheorem}

\begin{hbcorollary}{thm:p03-w03-p03-029-theory-form}{proved=annex:p03-w03-p03-029-theory-form}{Vaught's two-cardinal theorem for a countable theory}
Let \(\Lang\) be a countable first-order language, possibly with function and
constant symbols, containing the distinguished unary predicate \(U\), and let
\(T\) be an \(\Lang\)-theory. If \(T\) has a two-cardinal model
\(\mathcal M\) (\cref{def:p03-w03-p03-029-kappa-lambda-model}), then \(T\)
has an \((\aleph_1,\aleph_0)\)-model \(\mathcal N\) with
\(\mathcal N\equiv\mathcal M\).
\end{hbcorollary}

The direction of the transfer is fixed: from any two-cardinal model, whatever
its cardinals, to the single pair \((\aleph_1,\aleph_0)\). In the arrow
notation this is \((\kappa,\lambda)\to_{\aleph_1}(\aleph_1,\aleph_0)\) for
all \(\kappa>\lambda\geq\aleph_0\). Shelah records the gap-one case of the
result in the stronger compactness form
\((\lambda^+,\lambda)\to'_{\aleph_1}(\aleph_1,\aleph_0)\) and credits it to
Vaught with the Morley--Vaught paper as the reference
\cite[p.~4, item~(F)]{Shelah2004PairAlephN}. Barwise refers to the same result
as Vaught's two-cardinal theorem, dated 1964, and notes that it yields an
abstract completeness theorem for the logic \(\mathcal L(Q_1)\)
\cite[Chapter~I, \S1.5, pp.~8--9]{BarwiseFeferman1985}. The fixed direction
and the single target pair are what separate this theorem from the
L\"owenheim--Skolem theorems, from the far-apart transfer and from Chang's
transfer; \cref{def:p03-w03-p03-029-not-conflate} compares the five results
item by item.

\begin{hbmethod}{def:p03-w03-p03-029-mv-method}{How the Morley--Vaught proof works}
The proof on p.~55 of the source proceeds by countable homogeneous models and
an \(\omega_1\)-chain with a fixed predicate
\cite[p.~55, proof of Theorem~6.2]{MorleyVaught1962}. Every step is carried
out in \cref{sec:p03-w03-p03-029-mv-construction} and in the annex.
\begin{enumerate}
  \item The L\"owenheim--Skolem theorem gives \(\mathfrak B_1\prec\mathfrak
        A\) containing \(R_j^{\mathfrak A}\) with
        \(\overline{\mathfrak B_1}=\overline{R_j^{\mathfrak A}}\).
        Since \(\mathfrak A\) is larger, \(\mathfrak B_1\) is a proper
        elementary substructure with the same predicate.
  \item Display (7) of the source (\cref{thm:p03-w03-p03-029-mv-7}) gives
        \((\mathfrak B_2,C)\equiv(\mathfrak A,\abs{\mathfrak B_1})\) and an
        isomorphism of \(\mathfrak B_2\) onto \(\mathfrak B_2\mathbin{|} C\).
        Lemma~6.1(a) (\cref{thm:p03-w03-p03-029-mv-6-1}) then gives a
        countable homogeneous system \(\mathfrak A_1\) with a proper
        elementary substructure \(\mathfrak A_0\) such that
        \(\mathfrak A_0\cong\mathfrak A_1\) and
        \(R_j^{\mathfrak A_0}=R_j^{\mathfrak A_1}\), the predicate being
        countably infinite.
  \item By recursion on \(\xi<\omega_1\) one builds systems
        \(\mathfrak A_\xi\cong\mathfrak A_0\), each a proper elementary
        extension of its predecessors and all with the same predicate. At a
        limit stage the union is countable and homogeneous by
        Lemma~6.1(c), and Lemma~6.1(b) shows that it is again isomorphic to
        \(\mathfrak A_0\), so the recursion continues.
  \item The union of the whole chain has power \(\omega_1\), its
        predicate is the countable set \(R_j^{\mathfrak A_0}\), and it is
        elementarily equivalent to \(\mathfrak A\).
\end{enumerate}
The argument uses no indiscernible sequence and no
Ehrenfeucht--Mostowski model. The given large model is used only in the
first step. After that the argument lengthens a chain of countable models to
length \(\omega_1\).
\end{hbmethod}

\subsection{The Morley--Vaught construction}
\label{sec:p03-w03-p03-029-mv-construction}

This subsection supplies the proof of Theorem~6.2 in the notation of its
source, together with the earlier results of the source on which that proof
rests. These are the three parts of Lemma~6.1 and display (7) of the
source and, through them, its Lemma~1.2, Lemma~3.1(b), Theorem~2.3(b),
Corollary~2.4(d), Theorem~2.6 and Theorem~2.8\('\). The source proves the results of its Sections~2 and~3 for
systems of an arbitrary infinite power \(\kappa\), for classes satisfying
J\'onsson's conditions (III)--(VI\(_\kappa\))
\cite[pp.~40--41]{MorleyVaught1962}. The proof of Theorem~6.2 applies them
only to denumerable systems, whose finite subsystems play the role of the
subsystems of power less than \(\kappa\). Each of them is stated below in that
denumerable case, with the locator of the source's general statement, and is
proved in full in the annex. Lemma~1.2 is stated and proved in the source's
full generality. Where the proof given here takes a different route from the
source's, the entry says so.

\begin{hbdefinition}{def:p03-w03-p03-029-mv-notions}{Systems, reducts and homogeneity in Morley--Vaught notation}
The following notions and symbols are those of the source.
\begin{enumerate}
  \item A \emph{relational system} \(\mathfrak A=\langle
        A,R_i\rangle_{i\in I}\) consists of a non-empty set \(A\), its
        \emph{universe} \(\abs{\mathfrak A}\), and finitary relations
        \(R_i^{\mathfrak A}=R_i\) over \(A\); \(I=I_{\mathfrak A}\) is the
        \emph{index set}, and \(\overline{\mathfrak A}\) is the power of
        \(\abs{\mathfrak A}\) \cite[pp.~38--39]{MorleyVaught1962}. Two
        systems with the same index set and the same arity at each index are
        \emph{similar}. \emph{Denumerable} means of power \(\omega\).
        Throughout this subsection index sets are countable, as in the
        source from p.~54 on.
  \item For \(\emptyset\neq X\subseteq\abs{\mathfrak A}\),
        \(\mathfrak A\mathbin{|} X\) is the subsystem with universe \(X\) and
        the relations of \(\mathfrak A\) restricted to \(X\), and
        \(\mathfrak B\subseteq\mathfrak A\) means
        \(\mathfrak B=\mathfrak A\mathbin{|}\abs{\mathfrak B}\). The union of
        a chain of systems is formed by taking unions of universes and of
        relations \cite[p.~39]{MorleyVaught1962}.
  \item The \emph{\(J\)-reduct} \(\mathfrak A\rceil J\) is
        \(\langle\abs{\mathfrak A},R_i^{\mathfrak A}\rangle_{i\in
        I_{\mathfrak A}\cap J}\). For finitary relations \(S_k\) over
        \(\abs{\mathfrak A}\), \((\mathfrak A,S_k)_{k\in K}\) is the expanded
        system, and \((\mathfrak A,d_k)_{k\in K}\) is \(\mathfrak A\) with
        distinguished elements \(d_k\). For a class \(\mathscr N\),
        \(\mathscr N\rceil J=\{\mathfrak A\rceil J:\mathfrak
        A\in\mathscr N\}\) \cite[p.~39]{MorleyVaught1962}. In
        \((\mathfrak A,X,g)\) a subset \(X\) enters as a singulary relation
        and a function \(g\) on \(\abs{\mathfrak A}\) enters through its
        graph, a binary relation.
  \item \(\mathfrak A\equiv\mathfrak B\) means that the two similar systems
        satisfy the same first-order sentences. \(\mathfrak A\prec\mathfrak
        B\) (\(\mathfrak A\) is an \emph{elementary subsystem} of
        \(\mathfrak B\)) means \(\mathfrak A\subseteq\mathfrak B\) and
        \((\mathfrak A,a)_{a\in\abs{\mathfrak A}}\equiv(\mathfrak
        B,a)_{a\in\abs{\mathfrak A}}\); it is the relation written
        \(\preccurlyeq\) elsewhere in this handbook. A class \(\mathscr N\)
        is an \emph{elementary type}, \(\mathscr N\in\mathbf{ET}\), if
        \(\mathscr N=\{\mathfrak B:\mathfrak B\equiv\mathfrak A\}\) for some
        infinite \(\mathfrak A\); then \(I_{\mathscr N}\) is the common index
        set \cite[p.~40]{MorleyVaught1962}.
  \item If \(\varphi\) is a formula whose free variables are among
        \(v_0,\ldots,v_{m-1}\), \(m\) least, then \(\varphi^{\mathfrak A}\)
        is the \(m\)-ary relation that holds of \(x_0,\ldots,x_{m-1}\)
        exactly when \(\varphi\) is satisfied in \(\mathfrak A\) by that
        assignment, and \(\mathfrak A^*=(\mathfrak
        A,\varphi^{\mathfrak A})_{\varphi\in F}\), where \(F\) is the set of
        all formulas of the language of \(\mathfrak A\)
        \cite[p.~47]{MorleyVaught1962}.
  \item \(\mathscr S_\omega(\mathfrak A)\) is the set of finite subsystems of
        \(\mathfrak A\), and \(\mathscr I(\mathscr K)\) is the class of all
        systems isomorphic to a member of \(\mathscr K\)
        \cite[p.~39]{MorleyVaught1962}. Thus \(\mathscr I\mathscr
        S_\omega(\mathfrak A^*)\) is the class of isomorphs of finite
        subsystems of \(\mathfrak A^*\).
  \item A denumerable system \(\mathfrak A\) is \emph{homogeneous} if every
        isomorphism of a finite subsystem of \(\mathfrak A^*\) into
        \(\mathfrak A^*\) can be extended to an automorphism of
        \(\mathfrak A^*\). This is the source's definition (p.~48, through
        Definition~2.2(a) on p.~41) at power \(\omega\), where the
        subsystems of power less than the power of the system are the finite
        ones.
  \item For a function \(f\) on a finite non-empty
        \(X\subseteq\abs{\mathfrak A}\) into \(\abs{\mathfrak B}\), with
        \(\mathfrak A\) and \(\mathfrak B\) similar, call \(f\) an
        \emph{elementary map from \(\mathfrak A\) to \(\mathfrak B\)} if
        \((\mathfrak A,x)_{x\in X}\equiv(\mathfrak B,f(x))_{x\in X}\). The
        source has no name for this; \cref{thm:p03-w03-p03-029-mv-3-1}
        shows that these are its isomorphisms between finite subsystems of
        \(\mathfrak A^*\) and of \(\mathfrak B^*\).
  \item For singulary \(R_{j_1}^{\mathfrak A}\neq\emptyset\) and a function
        \(f_1\) from a set \(J\) into \(I_{\mathfrak A}\), the
        \emph{relativized reduct} \(\langle R_{j_1}^{\mathfrak
        A},R_{f_1(j)}^{\mathfrak A}\rangle_{j\in J}\) is the system with
        universe \(R_{j_1}^{\mathfrak A}\) whose \(j\)-th relation is the
        restriction of \(R_{f_1(j)}^{\mathfrak A}\) to \(R_{j_1}^{\mathfrak
        A}\) (pp.~39 and~51 of the source).
\end{enumerate}
\index{homogeneous system (Morley--Vaught)}\index{elementary type}
\end{hbdefinition}

\begin{hblemma}{thm:p03-w03-p03-029-mv-3-1}{proved=annex:p03-w03-p03-029-mv-3-1}{Finite subsystems of \(\mathfrak A^*\) (source Lemma 3.1(b))}
Let \(\mathfrak A\) and \(\mathfrak B\) be similar, let \(X\) be a finite
non-empty subset of \(\abs{\mathfrak A}\), and let \(f\) be a function on
\(X\) onto \(Y\subseteq\abs{\mathfrak B}\). Then
\(\mathfrak A^*\mathbin{|} X\cong_f\mathfrak B^*\mathbin{|} Y\) if and only
if \((\mathfrak A,x)_{x\in X}\equiv(\mathfrak B,f(x))_{x\in X}\)
\cite[p.~47, Lemma~3.1(b)]{MorleyVaught1962}. Moreover:
\begin{enumerate}
  \item the automorphisms of \(\mathfrak A^*\) are exactly the automorphisms
        of \(\mathfrak A\), and every elementary map is injective;
  \item a denumerable \(\mathfrak A\) is homogeneous if and only if every
        elementary map from \(\mathfrak A\) to \(\mathfrak A\) extends to an
        automorphism of \(\mathfrak A\);
  \item \(\mathscr I\mathscr S_\omega(\mathfrak A^*)\subseteq\mathscr
        I\mathscr S_\omega(\mathfrak B^*)\) if and only if for every finite
        non-empty \(X\subseteq\abs{\mathfrak A}\) there is an elementary map
        from \(\mathfrak A\) to \(\mathfrak B\) with domain \(X\).
\end{enumerate}
\end{hblemma}

The source states Lemma~3.1 and says that it is easily verified. Its
Theorem~3.3, which the source cites in proving Lemma~6.1, relates
homogeneity of \(\mathfrak A\) to homogeneity of \(\mathfrak A^*\); as
printed it carries the hypothesis \(\overline{\mathfrak A}=\kappa>\omega,
\overline{I}_{\mathfrak A}\) \cite[p.~48, Theorem~3.3]{MorleyVaught1962}. For
the denumerable systems used here, the passage between \(\mathfrak A\) and
\(\mathfrak A^*\) is made by the definition together with
\cref{thm:p03-w03-p03-029-mv-3-1}.

\begin{hblemma}{thm:p03-w03-p03-029-mv-1-2}{proved=annex:p03-w03-p03-029-mv-1-2}{Joint consistency (source Lemma 1.2)}
\begin{enumerate}
  \item[(a)] If \(\mathfrak D\equiv\mathfrak D'\), then for some
        \(\mathfrak D''\equiv\mathfrak D\) we have \(\mathfrak D,\mathfrak
        D'\in\mathscr I\mathscr S(\mathfrak D'')\).
  \item[(b)] Suppose \(\mathfrak D\equiv\mathfrak D'\); \(x_k\in\abs{\mathfrak
        D}\) for each \(k\in K\); and \(S_k\) is a finitary relation over
        \(\abs{\mathfrak D'}\) for each \(k\in K'\). Then there is a system
        \(\mathfrak D''\equiv\mathfrak D\) such that, for some \(y_k\)
        \((k\in K)\) and \(T_k\) \((k\in K')\),
        \((\mathfrak D,x_k)_{k\in K}\equiv(\mathfrak D'',y_k)_{k\in K}\) and
        \((\mathfrak D',S_k)_{k\in K'}\equiv(\mathfrak D'',T_k)_{k\in K'}\).
\end{enumerate}
Here \(\mathscr S(\mathfrak D'')\) is the set of all subsystems of
\(\mathfrak D''\), and the systems may carry distinguished elements
\cite[p.~40, Lemma~1.2]{MorleyVaught1962}. For index sets of arbitrary size
the proof uses the Axiom of Choice through the compactness theorem.
\end{hblemma}

The source attributes 1.2(a) to A.~Robinson and Fra\"iss\'e and 1.2(b) to
A.~Robinson \cite[p.~40, footnote~3]{MorleyVaught1962}.

\begin{hblemma}{thm:p03-w03-p03-029-mv-2-4d}{proved=annex:p03-w03-p03-029-mv-2-4d}{One-point criterion for homogeneity (source Corollary 2.4(d), denumerable case)}
A denumerable system \(\mathfrak A\) is homogeneous if (and only if) the
following holds: whenever \(f\) is an elementary map from \(\mathfrak A\) to
\(\mathfrak A\) with finite domain \(X\) and \(a\in\abs{\mathfrak A}\), there
is \(b\in\abs{\mathfrak A}\) such that \(f\cup\{(a,b)\}\) is an elementary
map from \(\mathfrak A\) to \(\mathfrak A\).
\end{hblemma}

The source's class parameter \(\mathscr M\) is a class of relational
systems, and \(\mathscr S^{\mathscr M}_\kappa(\mathfrak A)\) consists of
its members that are subsystems of \(\mathfrak A\) of power less than
\(\kappa\) \cite[p.~39]{MorleyVaught1962}. Corollary~2.4 assumes that
\(\mathscr M\) is closed under chain unions (condition V), that every
subset of power less than \(\kappa\) of a member is contained in a
subsystem in \(\mathscr M\) of power less than \(\kappa\) (condition
\(\mathrm{VI}_\kappa\)), and that \(\mathfrak A\in\mathscr M\) has power
\(\kappa\). Under these hypotheses its criterion says that any
isomorphism of \(\mathfrak B_1\in
\mathscr S^{\mathscr M}_\kappa(\mathfrak A)\) into \(\mathfrak A\) must
extend to each larger
\(\mathfrak B_2\in\mathscr S^{\mathscr M}_\kappa(\mathfrak A)\)
\cite[pp.~40--43, Corollary~2.4(d)]{MorleyVaught1962}.

For the application here take \(\mathscr M\) to be all systems similar
to \(\mathfrak A^*\) and \(\kappa=\omega\). Both class conditions hold,
and the indicated subsystems are exactly the finite nonempty ones.
Adding one element at a time gives the displayed criterion. Its proof is
the back-and-forth argument of Theorem~2.3, which the source describes as
a generalization of Cantor's arguments for the rational order
\cite[pp.~42--43]{MorleyVaught1962}.

\begin{hblemma}{thm:p03-w03-p03-029-mv-6-1}{proved=annex:p03-w03-p03-029-mv-6-1}{Morley--Vaught Lemma 6.1}
\begin{enumerate}
  \item[(a)] If \(\mathscr N\in\mathbf{ET}\) and \(J\subseteq I_{\mathscr
        N}\), then there is a denumerable, homogeneous system \(\mathfrak
        A\in\mathscr N\rceil J\).
  \item[(b)] If \(\mathfrak A\) and \(\mathfrak A'\) are homogeneous and
        denumerable, \(\mathfrak A\equiv\mathfrak A'\), and \(\mathscr
        I\mathscr S_\omega(\mathfrak A^*)=\mathscr I\mathscr
        S_\omega(\mathfrak A'^*)\), then \(\mathfrak A\cong\mathfrak A'\).
  \item[(c)] If \(\mathfrak A_0\prec\mathfrak A_1\prec\cdots\prec\mathfrak
        A_n\prec\cdots\) and each \(\mathfrak A_n\) is homogeneous and
        denumerable, then so is \(\bigcup(\mathfrak A_n/n\in\omega)\).
\end{enumerate}
\cite[p.~55, Lemma~6.1]{MorleyVaught1962}.
\end{hblemma}

The source proves (a) from 3.3, 1.2(b), 2.6 and 2.8\('\); (b) from 3.3 and
2.3(b); and (c) from 3.3 and 2.4(d) \cite[p.~55]{MorleyVaught1962}. The
annex proof of (a) carries out, at power \(\omega\), the recursion in the
proof of Theorem~2.6(a), with its auxiliary condition (3)
\cite[pp.~43--44]{MorleyVaught1962}, and keeps every stage inside
\(\mathscr N\) as Theorem~2.8\('\) requires \cite[pp.~45--46]{MorleyVaught1962};
the one-point step is obtained from Lemma~1.2(b), as the source does in
proving its Theorem~3.6 \cite[p.~50]{MorleyVaught1962}. The proof of (b) is
the back-and-forth argument of Theorem~2.3(b)
\cite[pp.~42--43]{MorleyVaught1962}. The source credits the special case of
2.3(b) stated in 6.1(b) to Craig \cite[p.~43, footnote~5]{MorleyVaught1962}.

\begin{hbtheorem}{thm:p03-w03-p03-029-mv-7}{proved=annex:p03-w03-p03-029-mv-7}{Display (7) of the source, countable index set}
Let \(\mathfrak A\) have a countable index set, let \(j_1,j_2\in
I_{\mathfrak A}\) with \(R_{j_1}^{\mathfrak A}\) and \(R_{j_2}^{\mathfrak
A}\) singulary and non-empty, and let \(J\) be a countable set and
\(f_1,f_2\in I_{\mathfrak A}{}^J\) with
\(R_{f_1(j)}\) and \(R_{f_2(j)}\) of the same arity for each \(j\in J\). If
\[
  \langle R_{j_1}^{\mathfrak A},R_{f_1(j)}^{\mathfrak A}\rangle_{j\in J}
  \equiv
  \langle R_{j_2}^{\mathfrak A},R_{f_2(j)}^{\mathfrak A}\rangle_{j\in J},
\]
then there exists \(\mathfrak B\equiv\mathfrak A\) such that
\(\langle R_{j_1}^{\mathfrak B},R_{f_1(j)}^{\mathfrak B}\rangle_{j\in
J}\cong\langle R_{j_2}^{\mathfrak B},R_{f_2(j)}^{\mathfrak B}\rangle_{j\in
J}\).
\end{hbtheorem}

The source states (7) without the restriction to countable index sets
\(I_{\mathfrak A}\) and \(J\) \cite[p.~51, display~(7)]{MorleyVaught1962};
the restriction is the standing hypothesis under which Theorem~6.2 uses it,
and there \(J=I_{\mathfrak A}\). The case of uncountable index sets is not
proved here. The source proves (7), when
\(R_{j_1}^{\mathfrak A}\) is infinite, by taking a special system
\(\mathfrak B\equiv\mathfrak A\) (its Theorem~3.5) and observing through its
Theorem~3.7 that both relativized reducts are special of the same power, hence
isomorphic; when \(R_{j_1}^{\mathfrak A}\) is finite it takes
\(\mathfrak B=\mathfrak A\), citing Theorem~1.4 of Tarski's
\emph{Contributions to the theory of models} (its reference~21). It credits (7), in a
slightly less general form, to Specker \cite[p.~51 and
footnote~11]{MorleyVaught1962}. The annex proves the infinite case by a
different route: it builds a denumerable \(\mathfrak B\equiv\mathfrak A\) in
which both relativized reducts are homogeneous and realize the same finite
configurations, and then applies Lemma~6.1(b). The finite case is proved
directly. The source's route through special systems of uncountable power,
its Theorems~3.5 and~3.7, is not reproduced here.

According to the source, Vaught formulated (7) in this general form to
answer affirmatively a question of D.~Scott: whether some system
elementarily equivalent to \(\langle\omega,+,\cdot,<\rangle\) is isomorphic to
a proper initial segment of itself. The answer comes from applying (7) to a
proper elementary extension of that structure with the standard part
\(\omega\) as a singulary relation \cite[p.~51]{MorleyVaught1962}.

\subsection{What the theorem does not say}

\begin{hblemma}{thm:p03-w03-p03-029-gap-one}{proved=annex:p03-w03-p03-029-gap-one}{Reduction to a gap-one pair}
Assume the Axiom of Choice. Let \(\Lang\) be countable and let
\(\mathcal M\) be a \((\kappa,\lambda)\)-model with
\(\kappa>\lambda\geq\aleph_0\). There is \(\mathcal N\preccurlyeq\mathcal
M\) with \(U^{\mathcal N}=U^{\mathcal M}\) and
\(\abs{\dom\mathcal N}=\lambda^+\). In particular \(\mathcal N\) is a
\((\lambda^+,\lambda)\)-model of the complete theory of \(\mathcal M\), and if
\(\lambda=\aleph_0\) it is an \((\aleph_1,\aleph_0)\)-model.
\end{hblemma}

\begin{hbtheorem}{thm:p03-w03-p03-029-robinson}{proved=annex:p03-w03-p03-029-robinson}{Robinson's extensional example}
Let \(\Lang_E=\{U,E\}\) with \(U\) unary and \(E\) binary. For an infinite set
\(X\) let \(\mathcal R_X\) be the \(\Lang_E\)-structure with universe
\((X\times\{0\})\cup(\mathcal P(X)\times\{1\})\), with
\(U^{\mathcal R_X}=X\times\{0\}\), and with
\(E^{\mathcal R_X}=\{((x,0),(Y,1)):x\in Y\}\).
\begin{enumerate}
  \item Every \(\mathcal N\equiv\mathcal R_X\) satisfies
        \(\aleph_0\leq\abs{U^{\mathcal N}}\) and
        \(\abs{\dom\mathcal N}\leq2^{\abs{U^{\mathcal N}}}\).
  \item \(\mathcal R_\omega\) is a \((2^{\aleph_0},\aleph_0)\)-model, so
        its theory admits a two-cardinal pair.
  \item For every infinite cardinal \(\lambda\) and every
        \(\kappa>2^\lambda\), the theory of \(\mathcal R_X\) does not admit
        \((\kappa,\lambda)\).
  \item Every elementary extension of \(\mathcal R_\omega\) of cardinality
        greater than \(2^{\aleph_0}\) has an uncountable predicate.
  \item The theory of \(\mathcal R_\omega\) admits
        \((\aleph_1,\aleph_0)\).
\end{enumerate}
\end{hbtheorem}

Morley and Vaught credit this example to Raphael Robinson. They use it to
show that a theory with a two-cardinal model need not admit a prescribed
larger gap. They also remark that iterating the construction gives the
bounds \(\abs{\dom\mathcal N}\leq2^{2^{\abs{U^{\mathcal N}}}}\) and so on
\cite[p.~55]{MorleyVaught1962}. The theorem above proves only the first
level.

\begin{hbtheorem}{thm:p03-w03-p03-029-hypotheses}{proved=annex:p03-w03-p03-029-hypotheses}{Each hypothesis is needed}
Each of the following theories has models of all the listed types and has no
\((\aleph_1,\aleph_0)\)-model.
\begin{enumerate}
  \item \emph{Finite predicate.} In \(\{U\}\), the theory \(T_1\) stating
        that \(U\) has exactly one element and that there are at least
        \(n\) elements, for each \(n\geq1\), admits \((\kappa,1)\) for every
        infinite \(\kappa\).
  \item \emph{No gap.} In \(\{U\}\), the theory \(T_2\) stating
        \(\forall x\,Ux\) and that there are at least \(n\) elements, for
        each \(n\geq1\), admits \((\kappa,\kappa)\) for every infinite
        \(\kappa\).
  \item \emph{Uncountable language.} In \(\{U\}\cup\{c_\alpha:\alpha<\omega_2\}\)
        with constant symbols \(c_\alpha\), the theory \(T_3\) stating
        \(c_\alpha\neq c_\beta\) and \(\neg Uc_\alpha\) for all
        \(\alpha\neq\beta\), and that \(U\) has at least \(n\) elements for
        each \(n\geq1\), admits \((\aleph_2,\aleph_0)\).
\end{enumerate}
\end{hbtheorem}

\begin{hbwarning}{def:p03-w03-p03-029-not-conflate}{Vaught's theorem is not the other transfer theorems}
Five results are often run together under the heading of two-cardinal
transfer. Only the first is the theorem of this section.
\begin{enumerate}
  \item \emph{Vaught's \((\aleph_1,\aleph_0)\) theorem}
        (\cref{thm:p03-w03-p03-029-mv-6-2,thm:p03-w03-p03-029-theory-form}).
        Hypothesis: a countable language and some two-cardinal model of any
        size. Conclusion: one specific pair, \((\aleph_1,\aleph_0)\). No
        cardinal arithmetic is assumed. Proof: countable homogeneous models
        and an \(\omega_1\)-chain (\cref{def:p03-w03-p03-029-mv-method}).
  \item \emph{The downward L\"owenheim--Skolem theorem}
        (\cref{thm:downward-ls}). It bounds the universe of an elementary
        substructure and can keep a given set, such as the whole predicate,
        inside it. It never makes the predicate smaller. It therefore
        reduces any two-cardinal model to a model of type
        \((\lambda^+,\lambda)\)
        (\cref{thm:p03-w03-p03-029-gap-one}). The source makes the same
        observation: lowering \(\overline{\mathfrak A}\) with
        \(\overline{R_j^{\mathfrak A}}\) fixed is the obvious positive case
        \cite[p.~56]{MorleyVaught1962}. When \(\lambda=\aleph_0\) this
        already yields \((\aleph_1,\aleph_0)\), so Vaught's theorem has new
        content only when the predicate is uncountable.
  \item \emph{The upward L\"owenheim--Skolem theorem.} It enlarges the
        universe but gives no control of the predicate. In
        \cref{thm:p03-w03-p03-029-robinson}, every elementary extension of
        \(\mathcal R_\omega\) of size above \(2^{\aleph_0}\) has an
        uncountable predicate. There is no upward version of Vaught's
        theorem with the predicate held fixed.
  \item \emph{Vaught's theorem for cardinals far apart.} In Shelah's
        account, this is the result of Vaught's 1965 paper: if
        \(\lambda\geq\beth_\omega(\mu)\) and \(\lambda'>\mu'\), then
        \((\lambda,\mu)\to'_{\leq\mu'}(\lambda',\mu')\). Morley later gave
        another proof using the Erd\H{o}s--Rado theorem and indiscernibles
        \cite[p.~4, item~(G)]{Shelah2004PairAlephN}. The hypothesis is a
        gap of at least \(\beth_\omega\) above the predicate. The
        conclusion reaches every pair \((\lambda',\mu')\) with
        \(\lambda'>\mu'\), larger pairs included. The indiscernible method
        of \cref{thm:p03-em-models} belongs to this theorem, not to the
        \((\aleph_1,\aleph_0)\) theorem.
        \Cref{thm:p03-w03-p03-029-robinson} shows that the far-apart
        hypothesis cannot be replaced by a single gap: a gap of one
        exponential does not transfer to larger universes over the same
        predicate.
  \item \emph{Chang's gap-one transfer.} In Shelah's account, Chang proved
        that if \(\mu=\mu^{<\mu}\), then
        \((\lambda^+,\lambda)\to'_{\leq\mu}(\mu^+,\mu)\)
        \cite[p.~4, item~(H)]{Shelah2004PairAlephN}. This theorem moves
        upward to a successor pair, and its hypothesis is a
        cardinal-arithmetic condition on the target. Jensen's result under
        \(\mathbf V=\mathbf L\) covers singular \(\mu\), the case left
        open by Chang's hypothesis. By the
        \((\aleph_1,\aleph_0)\) theorem, the source pair \((\lambda^+,\lambda)\)
        in these results may be taken to be \((\aleph_1,\aleph_0)\)
        \cite[p.~5, item~(I) and the remark after (J)]{Shelah2004PairAlephN}.
        Vaught's theorem is used here only as a first step. When
        Morley and Vaught wrote, it was not known whether \(\omega_0\) and
        \(\omega_1\) in Theorem~6.2 could be replaced by \(\kappa\) and
        \(\kappa^+\) \cite[p.~56]{MorleyVaught1962}.
\end{enumerate}
The bibliography key \texttt{Vaught1965CardinalsFarApart} names the
far-apart paper of item~4. That paper is not the source of the
\((\aleph_1,\aleph_0)\) theorem.
\end{hbwarning}

\begin{hbexample}{def:p03-w03-p03-029-worked-example}{Two theorems applied to one structure}
Take \(\mathcal R_\omega\) from \cref{thm:p03-w03-p03-029-robinson}. Its
predicate \(\omega\times\{0\}\) is countably infinite and its universe has
size \(\aleph_0+2^{\aleph_0}=2^{\aleph_0}\), so it is a
\((2^{\aleph_0},\aleph_0)\)-model.
\begin{enumerate}
  \item Downward. \Cref{thm:p03-w03-p03-029-gap-one} with
        \(\lambda=\aleph_0\) gives an elementary substructure of size
        \(\aleph_1\) containing the whole predicate: the \(\aleph_0\)
        elements of the predicate together with \(\aleph_1\) elements
        outside it.
        This \((\aleph_1,\aleph_0)\)-model is the conclusion of Vaught's
        theorem, and the L\"owenheim--Skolem theorem alone produced it.
  \item Upward. A model of the same theory of size \(2^{2^{\aleph_0}}\)
        exists by the upward L\"owenheim--Skolem theorem. By item~1 of
        \cref{thm:p03-w03-p03-029-robinson}, its predicate has
        cardinality \(\nu\) with \(2^\nu\geq2^{2^{\aleph_0}}\), so
        \(\nu>\aleph_0\). No model of this theory has type
        \((2^{2^{\aleph_0}},\aleph_0)\).
\end{enumerate}
A two-cardinal model therefore guarantees the pair \((\aleph_1,\aleph_0)\),
but not every larger universe over a countable predicate.
\end{hbexample}

\emph{Relationships.} The structure of
\cref{thm:p03-w03-p03-007-two-cardinal-boundary}, with universe \(\omega_1\)
and predicate \(\omega\), is an \((\aleph_1,\aleph_0)\)-model in the sense of
\cref{def:p03-w03-p03-029-kappa-lambda-model}. Its complete theory is not
\(\aleph_1\)-categorical, and there it shows that the categoricity hypothesis
of \cref{thm:p03-w03-p03-007-definable-sets} cannot be dropped. The use of Vaught's theorem for the logic with the
quantifier \(Q_1\) is taken up with Keisler's completeness theorem
(\cref{thm:p03-keisler-q-one}).

\begin{proofboundary}{def:p03-w03-p03-029-boundary}{Proved and cited parts}
\Cref{thm:p03-w03-p03-029-mv-6-2} is proved in the annex along the
source's construction, together with Lemma~6.1(a)--(c), display (7) for
countable index sets, Lemma~1.2, and the denumerable cases of Lemma~3.1(b),
Theorem~2.3(b), Corollary~2.4(d), Theorem~2.6(a) and Theorem~2.8\('\) on
which they rest. The source's statements of Theorems~2.3, 2.6 and~2.8\('\)
and Corollary~2.4 for systems of uncountable power, its Theorem~3.3, and
display (7) for uncountable index sets are cited and not proved here. Display
(7) is proved by a route different from the source's, which goes through the
special systems of its Theorems~3.5 and~3.7. The passage from the relational source form to a countable language with
function and constant symbols, the gap-one reduction, Robinson's bound and
the necessity of each hypothesis are proved in full in the annex. The
far-apart and gap-one transfers in items~4 and~5 of
\cref{def:p03-w03-p03-029-not-conflate} are reported as Shelah states them.
They are neither proved nor used in any proof here.
\end{proofboundary}
